\section{Ameaças à Validade}

As ameaças foram organizadas pelas quatro validades de \citet{wohlin2012experimentation}. Esta seção apresenta as principais; a lista completa, com a mitigação adotada e a mitigação futura de cada ameaça, está no Apêndice~\ref{ap:ameacas}.

\textbf{Validade de conclusão.} As 30 execuções de uma configuração reavaliam os mesmos mutantes com entradas aleatórias diferentes: medem a variação da geração aleatória, e não a variação entre programas. Por isso, os $p$-valores devem ser lidos junto com o $\hat{A}_{12}$ e com o teste por mutante distinto, que usa os mutantes como unidade; todos os $p$-valores foram corrigidos pelo método de Holm. O tempo foi medido em uma máquina virtual compartilhada, e a ordem aleatória das configurações em cada rodada distribui as variações de carga entre elas.

\textbf{Validade interna.} A primeira ameaça é a equivalência entre as implementações, tratada pela especificação única escrita antes do código e pelos 11 cenários compartilhados aprovados nos dois lados. Diferenças próprias de cada linguagem permanecem: o Java devolve objetos que o PIT muta como valores de retorno, enquanto os status do JavaScript são textos que o StrykerJS muta; o recorte harmonizado existe para isolar esses efeitos. A segunda ameaça é a suíte forte ter sido escrita pelo mesmo pesquisador que conhece o código. Ela foi derivada regra a regra do \texttt{SPEC.md}, e os dois ajustes feitos após a rodada piloto estão declarados (Seção~\ref{sec:fluxo}). A terceira é a geração de entradas: o jqwik injeta casos de borda que dominam a amostra quando há poucas tentativas, o que faz parte do fator ecossistema mas limita a atribuição das diferenças à linguagem em si.

\textbf{Validade de construto.} O \ms{} mede a detecção de defeitos artificiais, não de defeitos reais. Como as ferramentas geram conjuntos de mutantes diferentes \citep{kintis2018tools,monteiro2022tools}, o fator ecossistema mistura linguagem, biblioteca de PBT e ferramenta de mutação; o recorte harmonizado reduz, mas não elimina, a parte da ferramenta. A classificação dos mutantes não detectados combina critérios objetivos com julgamento manual por leitura do código, confirmado pelo autor, sem cálculo de concordância com um segundo classificador independente; a classe e a justificativa de cada mutante estão versionadas no repositório.

\textbf{Validade externa.} A aplicação tem cerca de 300 linhas de domínio em cada linguagem, e as propriedades verificam o domínio diretamente, sem passar pela camada HTTP nem por persistência real. Os resultados descrevem o comportamento das ferramentas e dos tipos de propriedade nesse domínio e não devem ser generalizados para APIs grandes sem replicação.
